\section{Marco teórico y estado del arte}
\label{sec:theoretical-framework}

El transporte cooperativo enlaza decisiones que suelen estudiarse por separado: selección de robots, formación de la coalición, estimación de recursos, contacto y movimiento de la carga. Cada resultado conserva su dominio. Un Nash puede ser ineficiente; una suma de capacidades puede ocultar un torque irrealizable; el consenso no estabiliza la planta; una trayectoria sin colisión puede terminar bloqueada.

\subsection{Transporte cooperativo y contratos entre capas}
\label{subsec:tf-agv-transport}

Los AGV surgieron ligados a rutas e infraestructura predefinidas, mientras los AMR incorporaron percepción y planificación a bordo. En instalaciones actuales la diferencia es funcional, no comercial: una plataforma puede navegar con flexibilidad y seguir sometida a corredores, prioridades y recursos compartidos \parencite{leanh2006agv,azadeh2019warehouse}. La coordinación estudiada exige que cada robot revise su decisión con estado propio y mensajes vecinales.

Una carga excede a un solo robot por masa, dimensiones, torque o maniobrabilidad. La cooperación exige algo más que llegada simultánea. El equipo debe cerrar su composición, ocupar contactos compatibles, producir el \emph{wrench} requerido y conservar seguridad durante el transporte \parencite{an2023cooperativeReview}. La Figura~\ref{fig:tf-problem-gap} muestra esas dependencias. Las flechas no expresan una garantía conjunta: cada transición necesita un modelo y una prueba propios.

\begin{figure}[htbp]
  \centering
  \begin{tikzpicture}[
    font=\sffamily,
    stage/.style={rounded corners=2pt,draw=VIUGray!70,line width=0.75pt,
      minimum width=2.42cm,minimum height=1.55cm,text width=2.15cm,
      align=center,inner sep=3pt,font=\fontsize{8.2}{9.1}\selectfont},
    measured/.style={stage,draw=SPBlue,fill=SPBlue!7},
    decided/.style={stage,draw=SPPurple,fill=SPPurple!7},
    certified/.style={stage,draw=VIUOrange,fill=VIUOrange!8},
    applied/.style={stage,draw=SPGreen,fill=SPGreen!8},
    arrow/.style={-{Latex[length=2.1mm]},draw=VIUDark,line width=0.9pt},
    return/.style={-{Latex[length=2mm]},draw=SPRed,line width=0.75pt,dashed},
    band/.style={rounded corners=2pt,draw=SPPurple!75,fill=SPPurple!4,
      minimum width=13.25cm,minimum height=0.60cm,align=center,
      font=\fontsize{7.8}{8.6}\selectfont}
  ]
    \node[measured] (measure) at (0, 0) {\textbf{Medición local}\\pose, batería, contacto};
    \node[measured] (estimate) at (2.75, 0) {\textbf{Estimación}\\ocupación y recursos};
    \node[decided] (decision) at (5.50, 0) {\textbf{Coalición}\\preferencia y cierre};
    \node[certified] (certificate) at (8.25, 0) {\textbf{Certificado}\\capacidad y \emph{wrench}};
    \node[applied] (plant) at (11.00, 0) {\textbf{Planta}\\pose, seguridad y fallo};

    \draw[arrow] (measure) -- (estimate);
    \draw[arrow] (estimate) -- (decision);
    \draw[arrow] (decision) -- (certificate);
    \draw[arrow] (certificate) -- (plant);
    \draw[return] (plant.south) -- ++(0,-0.62) -| (measure.south);
    \node[font=\fontsize{7.4}{8.0}\selectfont,text=SPRed,fill=white,inner sep=1pt]
      at (5.50,-0.93) {déficit, pérdida de contacto, error de pose o retirada};
    \node[band] at (5.50,-1.55)
      {Red imperfecta y escala: radio, retardo, pérdidas, mensajes, CPU y memoria};
  \end{tikzpicture}
  \caption{Contratos que separan información, decisión, factibilidad física y ejecución.}
  \label{fig:tf-problem-gap}
  \viuownsource
\end{figure}

La modalidad primaria, Cargo, modela la carga y los robots acoplados como un cuerpo planar soportado. El empuje cooperativo usa contactos unilaterales y admite deslizamiento. El \emph{caging} añade una condición geométrica de confinamiento. Estas modalidades comparten la tarea, pero no el modelo de contacto ni la prueba de estabilidad.

\subsection{Asignación y formación de coaliciones}
\label{subsec:tf-mrta}

La asignación multi-robot de tareas (MRTA) clasifica relaciones entre robots, tareas y horizonte de decisión \parencite{gerkeyMataric2004MRTA,korsahStentzDias2013iTax}. El caso uno a uno admite asignación bipartita y el algoritmo húngaro \parencite{kuhn1955Hungarian}. Una carga que requiere varios robots introduce dependencias entre ganadores. Su formulación mínima incluye
\begin{equation}
  \sum_{k=1}^{K}x_{ik}\leq 1,\qquad
  \sum_{i=1}^{N}x_{ik}\geq n_k,\qquad
  \sum_{i=1}^{N}\boldsymbol c_i x_{ik}\succeq \boldsymbol r_k,
  \qquad x_{ik}\in\{0, 1\}.
  \label{eq:tf-coalition-resource-feasibility}
\end{equation}
La primera restricción impone exclusividad; las otras dos, cuota y recursos. Replicar una tarea en plazas independientes no preserva necesariamente simultaneidad ni compatibilidad. MILP y \emph{set partitioning} son oráculos adecuados para instancias pequeñas, aunque dependen de información global.

Las subastas distribuyen la adjudicación mediante precios. El algoritmo de Bertsekas resuelve la asignación clásica \parencite{bertsekas1988Auction}; MURDOCH usa contratos en equipos robóticos \parencite{gerkeyMataric2002MURDOCH}. CBBA construye paquetes locales y resuelve conflictos por consenso, y ACBBA admite mensajes fuera de orden \parencite{choiBrunetHow2009CBBA,johnson2011acbba}. Sus garantías no cubren, sin una extensión explícita, el contacto simultáneo de varios robots con una carga.

La formación de coaliciones toma al grupo como unidad. Los métodos \emph{anytime} de Shehory y Kraus, las coaliciones basadas en capacidades de Vig y Adams y los juegos hedónicos cubren composición y preferencia colectiva \parencite{shehoryKraus1998Coalition,vigAdams2006Coalition,dutta2021hedonic}. ASyMTRe combina esquemas sensoriales, motores y de comunicación; IQ-ASyMTRe descarta composiciones incompatibles con requisitos funcionales \parencite{parkerTang2006ASyMTRe,zhangParker2013IQASyMTRe}. Estas formulaciones mejoran la selección del equipo, pero una coalición estable o rica en recursos todavía puede ser mecánicamente inviable.

Los antecedentes más próximos de 2024--2025 estrechan la comparación. Shan
et al. resuelven mediante subasta descentralizada tareas de transporte colectivo
que llegan dinámicamente, con ventanas temporales y corrección de desviaciones
de agenda \parencite{shan2024collectiveTransport}. Dai et al. aprenden de forma
conjunta coaliciones y rutas para tareas de cardinalidad variable
\parencite{dai2024dynamicCoalition}; Bezerra et al. incorporan mapas espaciales,
intenciones y revisión local de asignación mediante MAPPO
\parencite{bezerra2025dynamicCoalition}. Qiu et al. actualizan coaliciones cuando
el consumo reduce la capacidad disponible \parencite{qiu2024payloadConsumption},
y Shida et al. aprenden a excluir temporalmente tareas inviables para evitar
bloqueos \parencite{shida2025infeasibleTasks}. Estos trabajos enlazan asignación
con tiempo, movimiento o revisión dinámica, pero usan capacidad nominal o
aprendida y no el certificado explícito de contacto, \emph{wrench}, seguridad y
recuperación exigido por el problema. Los métodos aprendidos constituyen
comparadores cercanos, no el método principal, porque su decisión no es
directamente interpretable.

\subsection{Información local, juegos y optimización distribuida}
\label{subsec:tf-games-consensus}

La comunicación se representa mediante $\mathcal G(t)=(\Robots,\mathcal E(t))$. En un grafo no dirigido y conectado, $\lambda_2(L)>0$ mide conectividad algebraica y condiciona la disipación del desacuerdo \parencite{bullo2009distributed}. Los grafos conmutados requieren propiedades temporales, como conectividad conjunta, y cotas sobre retardo y actualización \parencite{olfatiSaber2004switchingConsensus}. El consenso certifica acuerdo sobre una variable; no certifica el juego ni el transporte.

Los juegos poblacionales describen preferencias continuas sobre acciones. En un Nash poblacional, toda estrategia con masa positiva maximiza el \emph{payoff} bajo el estado agregado \parencite{sandholm2010population}. Los juegos de potencial añaden una función escalar cuya variación coincide con cada desviación unilateral:
\begin{equation}
  u_i(a_i',a_{-i})-u_i(a_i,a_{-i})
  =\Phi(a_i',a_{-i})-\Phi(a_i,a_{-i}).
  \label{eq:tf-exact-potential}
\end{equation}
Esta igualdad debe probarse para el \emph{payoff} implementado. Descuentos, proyecciones o información obsoleta pueden romperla. Replicator, Smith, BNN y logit tampoco son nombres intercambiables: cambian la comparación de incentivos, el soporte y los estados estacionarios \parencite{sandholm2010population}. La masa continua necesita un cierre entero antes de representar robots identificables.

Un equilibrio de Nash generalizado incorpora restricciones compartidas. Bajo convexidad y regularidad puede formularse como desigualdad variacional; el equilibrio variacional comparte multiplicadores entre agentes \parencite{yi2019operatorGNE,franci2022stochasticGNE}. Los métodos primal--dual actualizan decisiones y multiplicadores, con garantías condicionadas por monotonicidad, conectividad y ganancias \parencite{cherukuri2016primalDual,koshal2016aggregative}. Ninguna de estas propiedades convierte por sí sola el equilibrio en óptimo social.

La pasividad y Lyapunov permiten estudiar la interconexión con la planta. Una prueba requiere variable de almacenamiento, suministro, dominio y signo de la derivada o diferencia \parencite{vanderschaf2017passivity,ortega2002passivity}. La etiqueta \emph{pasivo} no basta si las variables de puerto o los modelos de contacto cambian entre bloques.

\subsection{Contacto, formación y \emph{wrench}}
\label{subsec:tf-rigidity-wrench}

La conectividad de una formación no implica rigidez. La rigidez geométrica determina si las restricciones relativas fijan localmente la forma, salvo traslación y rotación \parencite{desai2001formation,bullo2009distributed}. En Cargo también importan el soporte vertical y los límites de actuador.

Si $\boldsymbol\lambda_C$ agrupa esfuerzos de contacto y $G_C(q)$ es la matriz de agarre, el \emph{wrench} aplicado es
\begin{equation}
  W_C=G_C(q)\boldsymbol\lambda_C,
  \qquad \boldsymbol\lambda_C\in\mathcal U_C.
  \label{eq:tf-grasp-map}
\end{equation}
La literatura distingue movimiento, fuerzas resultantes y fuerzas internas \parencite{nakamura1989dynamics,yoshikawa1993coordinated}. La geometría y el tipo de contacto determinan las propiedades de cierre \parencite{bicchi1995closure}. Por eso una suma escalar de capacidades no sustituye la pertenencia $W_k^{\mathrm{req}}\in\mathcal W_C(q)$, con $\mathcal W_C(q)=\{G_C(q)\boldsymbol\lambda:\boldsymbol\lambda\in\mathcal U_C\}$.

En empuje, $\mathcal U_C$ impone unilateralidad y fricción. El \emph{caging} exige que la configuración de robots confine el objeto; rodearlo no basta \parencite{pereira2004caging,fink2008caging}. Los métodos de transporte restringido pueden repartir fuerzas y seguir la pose para un equipo dado \parencite{alonsomora2017transport}. El problema adicional consiste en seleccionar ese equipo con información local.

\subsection{Seguridad, fallos y red imperfecta}
\label{subsec:tf-industrial-constraints}

ORCA construye velocidades recíprocamente admisibles bajo un modelo geométrico \parencite{vandenberg2011orca}. Las funciones de barrera restringen el control para mantener un conjunto seguro cuando el estado inicial, el modelo y el problema de control satisfacen sus supuestos \parencite{ames2017cbf,wang2017barrier}. Ambos enfoques pueden detenerse en congestión. Seguridad y vivacidad son propiedades distintas.

Los fallos exigen revisar la coalición y el contacto. ALLIANCE usa impaciencia y aquiescencia para activar conductas alternativas, mientras las subastas dinámicas reasignan tareas después de eventos \parencite{parker1998alliance,nanjanath2006dynamicAuctions}. En una carga rígida, retirar un robot cambia además el soporte y el conjunto de \emph{wrench}; la reparación lógica necesita un nuevo certificado físico.

Retardo, pérdida y topología alteran la información disponible. El consenso con enlaces imperfectos depende del modelo estadístico y de la conectividad \parencite{kar2009imperfectConsensus,nedic2010constrainedConsensus}; ACBBA también se ha estudiado con pérdida de mensajes \parencite{rantanen2018lossyAcbba}. Estas garantías no cubren particiones persistentes ni demuestran que un equilibrio visible sea globalmente consistente.

\subsection{Comparación y brecha de investigación}
\label{subsec:tf-synthesis-gap}

El contrato básico de cada familia se compara en la Tabla~\ref{tab:tf-comparison}. ``Parcial'' indica que la propiedad depende de un módulo adicional o de supuestos ajenos al problema tratado.

\begin{table}[!htbp]
  \centering
  \caption{Cobertura de las familias metodológicas relevantes.}
  \label{tab:tf-comparison}
  \scriptsize
  \setlength{\tabcolsep}{2.4pt}
  \begin{tabularx}{\textwidth}{@{}p{2.0cm}p{1.25cm}p{1.35cm}p{1.45cm}X p{2.2cm}@{}}
    \toprule
    Familia & Coalición & Física & Información local & Garantía típica & Límite para este problema \\
    \midrule
    Asignación/MILP & Sí & Parcial & No & óptimo o cota del optimizador & coste global y combinatorio \\
    Subasta/CBBA & Parcial & No & Sí & asignación bajo supuestos de puntuación y red & no certifica contacto \\
    Juego de coalición & Sí & Parcial & Sí & estabilidad o mejora local & no implica óptimo social \\
    Juego poblacional & Relajada & No & Parcial & convergencia bajo estructura del juego & exige cierre entero \\
    Formación/\emph{wrench} & Equipo dado & Sí & Parcial & rigidez, reparto o seguimiento & no selecciona el equipo \\
    ORCA/CBF & No & Parcial & Sí & evitación o invariancia condicionada & puede bloquearse \\
    MARL/GNN & Aprendida & Parcial & Sí & rendimiento empírico & interpretación y garantías limitadas \\
    \bottomrule
  \end{tabularx}
  \viuownsource
\end{table}

La brecha no es la falta de métodos aislados. MRTA cierra asignaciones; el consenso alinea estimaciones; un juego regula incentivos; el control mueve un equipo dado. Falta comprobar que las interfaces entre esos resultados conservan exclusividad, factibilidad mecánica, seguridad, recuperación y calidad bajo red imperfecta. En el corpus revisado no se identificó una solución única que cubra esas condiciones con un mecanismo distribuido y explicable. La afirmación se limita a las fuentes verificadas y no pretende demostrar inexistencia universal.

La evaluación añade las condiciones una a una y conserva los fallos que impiden extender cada conclusión al sistema completo.
